\section{引言}
\label{sec:intro}

\subsection{经验问题：重尾的、不对称的信息扩散}
在线社交网络中的信息扩散是显著\emph{重尾}的：极少数级联会触达巨大的受众，而
绝大多数级联迅速消亡 \citep{barabasi2005,goel2012,leskovec2007}。更关键的
是，经验研究已经确立了这种扩散的\emph{不对称性}：\citet{vosoughi2018} 对十余
年间推特上约 12.6 万条信息级联的分析表明，虚假内容在各个信息类别中都比真实
内容传播得更快、更深、更广，而新颖性与情绪反应是主要驱动因素。

这一不对称性目前缺少机制层面的解释。主流做法是把它归因于内容属性（耸动性、
新颖性）与受众心理（新奇偏好、情绪唤起），这些当然是重要的，但它们把"公众
易感性"当作常量，而非一个随环境变化的\emph{状态变量}。本文的核心问题是：
\emph{公众对误信息的易感性是否由一个共振机制所支配？}换言之，是否存在某个
特定的"社会噪声"水平，使得群体对阈下信号——无论真伪——的响应被最大化？

重尾性本身具有深刻的动力学含义。在重尾环境中，罕见的大事件（病毒式帖子、
意见领袖表态、协调化放大、机器人账号的同步推送）主导系统的涨落，中心极限定理
不再适用，基于高斯噪声的模型会系统性地低估极端事件的频率与后果。这提示我们
用 $\alpha$-稳定分布——重尾涨落的自然建模语言 \citep{applebaum2009}——来刻画
社会噪声，而不是沿用高斯假设。

\subsection{意见动力学与社会物理学：机制性谱系}
社会物理学为意见与传播现象提供了丰富的机制性模型库。\citet{granovetter1978}
的阈值模型确立了"个体门槛的异质性会导致集体结果的剧烈不连续变化"这一基本洞见；
\citet{galam2002} 表明不灵活的少数派（顽固者）可以推翻多数；\citet{watts2002}
给出了随机网络上全局级联的一般条件；\citet{kempe2003} 把影响最大化表述为一个
组合优化问题；\citet{castellano2009} 对这一脉络做了系统综述。

然而，这条脉络的主流模型是\emph{确定性}或\emph{简单随机}的：噪声通常只作为
扰动项进入，从未被当作具有\emph{功能}的要素。与此同时，传播研究
\citep{goel2012,leskovec2007} 主要关注结构与可达性，较少关注\emph{时间上的
相干性}——即群体是否\emph{同步地}响应。本文正是要补上这一维度：噪声不只是
需要被平均掉的扰动，它决定了群体能否响应。

\subsection{随机共振：一种依赖于噪声的易感性机制}
随机共振（SR）指的是一个反直觉的现象：非线性系统对一个微弱的、通常处于阈值
以下的输入信号的响应，会被\emph{中等强度}的噪声所\emph{优化}。这一概念由
\citet{benzi1981,benzi1982} 提出用以解释地球冰期的重现，此后在
\citet{mcnamara1989} 的两态约化与 \citet{gammaitoni1998} 的综述中得到系统
表述，并在物理、生物与神经科学中被广泛记录。

SR 之所以与误信息问题高度相关，是因为它给出了一条\emph{关于噪声功能的一般
原理}：在具有阈值的系统中，噪声不是对信号单纯的污染，它可以是信号得以被
传递的\emph{必要条件}。典型的 SR 曲线呈倒 U 形：噪声太小，系统被"冻结"在
初态，输出中不含输入信息；噪声太大，输入--输出相关性被噪声破坏；只有在中等
噪声处，噪声诱导的跃迁速率与信号时标相匹配，宏观响应达到最大。

把这一机制用于公共领域，就得到一个可检验的、反直觉的命题：\emph{存在一个
特定的社会噪声水平，使得公众最容易被（任何）阈下信号所动员}。本文的任务
就是把这一命题形式化、证明它、并用仿真检验它。

我们采用三类互补的指标，以避免结论依赖于单一的度量选择：(i) 经典的
\emph{信噪比}；(ii) 由非微扰矩展开得到的\emph{谱放大因子}
$\eta=2|M_1|/A$——即动力学磁化率在驱动频率处的模 \citep{liu2018levy}，它由于
被驱动幅值归一化而可以在不同噪声统计之间公平比较；(iii) 针对\emph{非周期}
（随机二值）阈下信号的输入--输出\emph{互信息} $I(S;Y)$，它完全不预设周期性
\citep{kang2020aperiodic,patel2008,kosko2009}。三个先验不同的量若在同一处
同时呈现峰值，就排除了"峰只是某种噪声本底估计方式的人工产物"这一解释。

\subsection{研究意义}
\label{sec:significance}
本文的研究意义可以从三个层面陈述。

\paragraph{(i) 理论层面：为误信息的易感性提供一个可计算的机制模型。}
现有误信息研究大多是描述性/统计性的：测量传播深度、速度、广度
\citep{vosoughi2018}，或做基于主体的传播模拟 \citep{goel2012}。这类研究能
告诉我们"发生了什么"，但难以回答"为什么在这个条件下发生、换一个条件会怎样"。
本文把公众易感性归结为一条可计算的共振曲线，并给出\emph{无参数}的定量预言
（最优噪声强度 $D^{*}=\Delta U/2$），从而使"易感性"成为一个可预测、可干预的
量，而非一个事后描述的标签。

\paragraph{(ii) 方法层面：把重尾噪声与空间传播同时引入意见动力学。}
经典 SR 理论几乎全部是高斯且无空间结构的；而经典意见动力学又几乎不考虑
重尾噪声的功能性作用。本文首次把"耦合 + 重尾 + 空间传播 + 信息论指标 +
误信息语义"放在同一个模型中处理，并给出严格的适定性基础。

\paragraph{(iii) 实践层面：给出可操作的治理含义。}
本文的结果直接支持三条治理建议（第~\ref{sec:governance} 节）：同步性应被
当作风险指标而非健康指标；干预应部署在群体最可激发的时段；治理应靶向
波前而非仅源头。其中每一条都与当前实践有实质差异。

\subsection{国内外研究现状与本文定位}
\label{sec:literature}
相关研究可梳理为四条相对独立的脉络，本文位于其交叉点上目前仍然缺失的位置。

\paragraph{脉络一：随机共振的经典与推广理论。}
自 \citet{benzi1981} 以来，SR 理论在 \citet{mcnamara1989} 与
\citet{gammaitoni1998} 中系统化，并推广到非周期信号、超阈值信号、阵列与
网络的协同共振。但这条脉络的工作几乎全部是\emph{高斯}的，且大多无空间结构。

\paragraph{脉络二：\Levy\ 噪声与分数阶动力学。}
\citet{jespersen1999}、\citet{chechkin2003} 建立了 \Levy\ 驱动系统的分数阶
Fokker--Planck 形式体系；\citet{applebaum2009} 给出严格随机分析基础。在这一
脉络中 SR 被研究过，但通常限于\emph{单体}系统 \citep{liu2018levy}。

\paragraph{脉络三：意见动力学与社会物理学。}
从 \citet{granovetter1978} 的阈值模型到 \citet{galam2002} 的顽固者模型，
再到 \citet{castellano2009} 的综述；传播方面则有 \citet{watts2002}、
\citet{kempe2003}、\citet{leskovec2007}。这条脉络的主流是确定性或简单随机的。

\paragraph{脉络四：误信息的经验研究。}
\citet{vosoughi2018} 的经典研究确立了虚假内容的传播优势；\citet{goel2012}
刻画了在线扩散网络的结构。但这一脉络基本是描述性的，缺乏可推广的机制模型。

\paragraph{本文定位与贡献。}
四条脉络各自成熟，但交叉处是空白：\emph{在带空间耦合的意见动力学中，重尾
噪声能否产生随机共振？若能，经典理论预言的峰位是否仍成立？这对误信息治理
意味着什么？}具体而言，本文贡献是：
\begin{enumerate}[leftmargin=2em,itemsep=2pt]
  \item 建立模型的适定性与矩估计，并证明 $D\to0$ 时的均方收敛，为"小噪声端
        响应必然消失"提供严格基础；
  \item 导出空间分数阶 Fokker--Planck 方程与非微扰矩层级，得到谱放大因子；
  \item 给出平均场约化与有序--无序相变，把共振定位在相变线附近；
  \item 由 Kramers 逃逸率给出无参数峰位预言 $D^{*}=\Delta U/2=0.125$；
  \item 证明阈下二值输入下互信息 $I(S;Y)$ 的非单调性（非周期 SR 的存在性）；
  \item 数值上记录重尾驱动下共振峰的右移与展宽，并用首次穿越时间测量独立地
        确认其机制为跳跃介导的势垒跨越；
  \item 识别"误信息脆弱性窗口"，刻画辟谣的非线性回报，并给出平台治理含义。
\end{enumerate}

